\chapter{Information Gathering}
\label{ch:info_gathering}

Information gathering is the first phase of any penetration test. Before touching
the target, I need to understand what I'm dealing with --- domains, IPs, services,
infrastructure, and any publicly available information that reduces my attack surface
from unknown to known.

\begin{notebox}[Course emphasis]
Always start with passive reconnaissance before moving to active reconnaissance.
This is not just good practice --- it is the professional and legal way to approach
an engagement.
\end{notebox}

% ─────────────────────────────────────────────────────────────────────────────
\section{Target Scoping}
\label{sec:target_scoping}

Before any testing begins, the client defines what is in scope. This is a legal
and professional boundary.

\textbf{What scope defines:}
\begin{itemize}
  \item A single domain (e.g., \ic{example.com}) or specific subdomains
  \item One or more IP addresses or a full network range (e.g., \ic{192.168.1.0/24})
  \item Specific services or applications to test
  \item Time windows allowed for testing
\end{itemize}

\begin{warnbox}[Stay in scope]
Testing outside the defined scope is illegal, even accidentally. Always confirm
what you are allowed to touch before running any tool.
\end{warnbox}

\begin{takebox}
\textbf{Key takeaway:} Scope definition saves you from legal liability and stops you
from wasting time collecting data about systems you are not allowed to test.
A client giving you \ic{192.168.1.0/24} does not mean you can touch \ic{10.0.0.1}.
\end{takebox}

% ─────────────────────────────────────────────────────────────────────────────
\section{Passive vs Active Reconnaissance}
\label{sec:passive_active_recon}

\subsection{Passive Reconnaissance}

Passive recon means collecting information without directly interacting with the
target systems. No packets go to the target. This keeps you invisible.

\textbf{Characteristics:}
\begin{itemize}
  \item No direct connection to target servers
  \item Very low risk of detection
  \item Always performed first
\end{itemize}

\textbf{What you can find passively:}
\begin{itemize}
  \item Domain registration information (WHOIS)
  \item DNS records
  \item Public website content
  \item Search engine results (Google dorking)
  \item Publicly available email addresses
  \item Historical snapshots (Wayback Machine)
  \item Data breach databases (HaveIBeenPwned)
\end{itemize}

\subsection{Active Reconnaissance}

Active recon means sending traffic directly to the target. This is detectable. You
do this only after passive recon has given you a clear picture of the attack surface.

\textbf{Characteristics:}
\begin{itemize}
  \item Sends traffic to the target --- logs are created
  \item Increased risk of detection
  \item Typically performed after passive recon
\end{itemize}

\textbf{What you discover actively:}
\begin{itemize}
  \item Live hosts on a network
  \item Open ports and running services
  \item Service versions and banners
  \item Network topology and responses
\end{itemize}

% ─────────────────────────────────────────────────────────────────────────────
\section{Recon Strategy --- The Four-Step Approach}
\label{sec:recon_strategy}

I follow this every single engagement. It keeps me structured and stops me
jumping into exploitation before I actually understand the target.

\begin{enumerate}
  \item \textbf{Define the Target} --- identify the domain, IP, or network range.
        Confirm what is in scope.
  \item \textbf{Perform Passive Reconnaissance} --- gather public information,
        identify potential attack surfaces without touching the target.
  \item \textbf{Perform Active Reconnaissance} --- scan for live hosts, open ports,
        running services, and version information.
  \item \textbf{Document Everything} --- record all findings before moving to
        enumeration. What you skip here costs you later.
\end{enumerate}

% ─────────────────────────────────────────────────────────────────────────────
\section{Passive Recon Tools}
\label{sec:passive_tools}

\subsection{WHOIS --- Domain Registration Lookup}

WHOIS returns registration data for a domain: registrar, owner contact details,
name servers, and registration dates.

\begin{whybox}[Why this matters]
WHOIS can reveal the registrant's email, which can be used for OSINT, phishing
context, or identifying admin contacts. Name server info helps with DNS enumeration.
\end{whybox}







\begin{termbox}
whois target.com
\end{termbox}

\begin{takebox}

\begin{whybox}[Why this matters]

\begin{notebox}[Course emphasis]

\screenshot{assets/images/Screenshot_2025-07-10_at_11.36.07_PM.png}{WHOIS lookup showing domain registration information}

\subsection{host --- DNS Lookup}

The \tool{host} command resolves a domain to its IP addresses and can retrieve
different record types.

\begin{termbox}
host target.com            # A record (IPv4)
host -t mx target.com      # Mail exchange records
host -t ns target.com      # Name server records
host -t txt target.com     # TXT records (SPF, DKIM, etc.)
\end{termbox}

\subsection{HaveIBeenPwned --- Breach Data}

\url{https://haveibeenpwned.com} checks whether email addresses or domains appear
in publicly known data breaches. Useful for understanding whether credentials for
the target organisation may be circulating.

\begin{notebox}
If a company's email domain appears in breach data, assume some employee passwords
are compromised. This is a starting point for credential stuffing attacks.
\end{notebox}

\subsection{Netcraft --- Website Footprinting}

Netcraft (\url{https://sitereport.netcraft.com}) reveals:
\begin{itemize}
  \item Web server technology and version
  \item Hosting provider and IP history
  \item SSL certificate information
  \item Risk rating
\end{itemize}

\subsection{DNS Recon}
\label{subsec:dns_recon}

\tool{dnsrecon} automates DNS enumeration --- it queries multiple record types and
can perform zone transfer attempts.

\begin{termbox}
dnsrecon -d target.com           # Standard recon on domain
dnsrecon -d target.com -t std    # Standard enumeration
dnsrecon -d target.com -t zt     # Zone transfer attempt
\end{termbox}

\screenshot{assets/images/Screenshot_2025-07-11_at_9.35.21_PM.png}{dnsrecon output showing DNS records for a target domain}

\subsection{DNS Zone Transfer}
\label{subsec:zone_transfer}

A zone transfer (AXFR) is a DNS server mechanism that replicates DNS records to
another server. If a DNS server is misconfigured to allow zone transfers from any
IP, an attacker can retrieve the entire DNS zone --- all subdomains, IPs, and records.

\begin{whybox}[Why attackers love this]
A successful zone transfer hands you a complete map of the target's internal DNS
infrastructure. Every subdomain, internal host, and IP in one response.
\end{whybox}

\begin{termbox}
# Using dig to attempt a zone transfer
dig axfr @ns1.target.com target.com

# If internal nameserver is known
dig axfr @internal_nameserver internal.target.com
\end{termbox}

\screenshot{assets/images/Screenshot_2025-07-17_at_5.03.20_PM.png}{DNS zone transfer attempt showing returned records}

\subsection{Subdomain Enumeration with Sublist3r}

\tool{Sublist3r} searches public sources for subdomains: search engines, certificate
transparency logs, DNS databases.

\begin{termbox}
sublist3r -d target.com
sublist3r -d target.com -o output.txt    # Save results
\end{termbox}

\screenshot{assets/images/Screenshot_2025-07-17_at_5.00.48_PM.png}{Sublist3r discovering subdomains for a target}

% ─────────────────────────────────────────────────────────────────────────────
\section{Learning Output}
\label{sec:info_gathering_learning}

\begin{takebox}
\textbf{What I Learned}
\begin{itemize}
  \item Passive recon is more powerful than it looks --- WHOIS and DNS alone can
        reveal internal infrastructure if zone transfers are misconfigured.
  \item The order matters: passive first, active second. Jumping straight to nmap
        before WHOIS means you might miss context that changes your whole approach.
  \item Sublist3r + DNS recon together give a much fuller picture than either alone.
\end{itemize}

\textbf{Key Takeaways}
\begin{itemize}
  \item Always check WHOIS, DNS records, and Netcraft before touching the target.
  \item Zone transfers are a quick win --- try them on every nameserver found.
  \item HaveIBeenPwned is useful for understanding credential risk before you even
        run a single scan.
\end{itemize}

\textbf{Enumeration Mindset --- What to Check First}
\begin{enumerate}
  \item WHOIS the domain --- who owns it, what name servers
  \item \ic{host -t ns target.com} --- get the authoritative name servers
  \item Attempt zone transfer on each name server
  \item Run Sublist3r for subdomain discovery
  \item Check Netcraft for server technology and hosting history
\end{enumerate}

\textbf{Common Mistakes}
\begin{itemize}
  \item Skipping passive recon and going straight to nmap --- noisy and unnecessary
  \item Not documenting findings before moving on --- you will forget IP-to-hostname
        mappings and regret it during exploitation
  \item Assuming zone transfers always fail --- they succeed more often than expected
        on misconfigured legacy infrastructure
\end{itemize}

\textbf{Real-World Impact}
\begin{itemize}
  \item If zone transfer succeeds: full internal network map, no guessing needed
  \item Breach data exposure: existing credentials may work on the target
  \item Subdomain discovery often reveals staging, dev, or admin portals not visible
        from the main domain
\end{itemize}
\end{takebox}

\end{notebox}
\end{whybox}
\end{takebox}